\par\Needspace{6\baselineskip}\medskip
\begin{center}\hypertarget{native-177-article130}{}\pdfbookmark[2]{130.}{article-130}130.\end{center}
The ten methods explained from Article 124 onward rest on the assumption that approximate
values are already available for the celestial body's distances from Earth or for the orientation of
the orbital plane. This requirement clearly presents no difficulty when the task is to use more
widely separated observations to correct orbital dimensions whose approximate values are already
known, for example, from an earlier calculation based on other observations. The question remains
how to begin the first calculation when all the dimensions of the orbit are still entirely unknown.
This is by far the most important and difficult case of our problem, as one might expect from the
analogous problem in the theory of comets. It is well known how long that problem vexed
mathematicians and how many fruitless attempts it prompted. If the solution follows the procedure
explained from Article 119 onward, it must clearly satisfy the following conditions for our problem
to count as correctly solved. \emph{First}, the quantities $x,y$ must be chosen so that approximate
values for them can be obtained from the very nature of the problem, at least as long as the
heliocentric motion of the celestial body between observations is not too large. \emph{Second},
small variations in the quantities $x,y$ must produce variations that are not too large in the
quantities derived from them. Otherwise, any errors in the assumed values of the former might
prevent us from treating the latter as approximations. \emph{Third} and finally, we require that the
operations by which we proceed step by step from the quantities $x,y$ to $X,Y$ should not be too
lengthy.

These conditions provide a criterion for judging the merits of any method, which repeated use will
reveal still more clearly. The method we are now about to explain, which may in a sense be regarded
as the most important part of this work, satisfies these conditions so well that it seems to
\SourceLeaf{178}leave nothing further to be desired. Before explaining it in the form most convenient for
practical use, we will give some preliminary considerations to clarify and open up an approach that
might otherwise seem somewhat obscure and less readily apparent.
